%% ============================================================================
%% II. Related Work                                           target: ~0.5 page
%% Stage 1: cite only the work most directly tied to the central problem (what
%% is needed to establish novelty). Stage 2 requires a complete section.
%% Anonymity: cite the authors' own prior work in the third person.
%% RSS 2027 checks references automatically and desk-rejects hallucinated
%% citations, so check every entry you add to references.bib.
%% Merged with the notes in material/2_Related_work.tex (earlier document).
%% ============================================================================
\section{Related Work}
\label{sec:related}



Manipulation of DLOs has applications in a wide range of human activities \cite{zhu_challenges_2021}, and has been a very active field of robotics  research in recent years \cite{sanchez_robotic_2018,caporali2026robotic}. A large,  body of work has been dedicated to estimating \cite{10167643,11248900,11568531} or tracking \cite{9928322,dinkel_dlo-splatting_2025,caporali2025robotic} the state of DLOs as a prerequisite for manipulations tasks. However, this work relies heavily on cameras which struggle with capturing the thin DLOs in the presence of clutter or when occlusions occur \cite{cop2021new}. For this reason, some recent works combine visual and haptic sensors when working with DLOs~\cite{nadon_multi-modal_2018,pecyna_visual-tactile_2022,11246285,11482717,caporali2026vision}. 

It is well-established that haptic information can provide quality information regarding the state of grasped DLOs \cite{bretl2014quasistatic,ms-lcwrtsemp-23,monguzzi2024force} and is even sufficient for manipulation under certain constraints \cite{smbb-mpdamer-20,slk-ftmdloeffi-22}. This observation is especially important for insertion and threading tasks where grasped objects, not necessarily DLOs, are routed through constrained spaces where they are partially or fully obstructed from cameras by the environment and the robot. Haptic information has been successfully used in the non-deformable object insertion task to complement camera input \cite{caporali2026vision} or as the main sensing modality \cite{kamijo2026tactile}. When DLOs are involved, contact of the object with the environment carries information about the state of the object as well as its alignment, thus increasing the significance of haptic sensors. Tactile sensors are used in \cite{de2018integration} and \cite{zanella2019dlo} to estimate the DLO's mechanical properties for the insertion process. In \cite{lvw-crcrisb-24}, A F/T sensor is used to guide a DLO through short tunnels whose shape is unknown. Another work introduces several DLO routing primitives, including insertion, that utilize haptic sensors \cite{wjly-cratmp-22}.   

Contact-aided motion planning is mostly useful for soft continuum and soft robots that can use the environment to change their orientation \cite{marchese2016dynamics,rao2024towards}. Due to the wide use of continuum robots in medical settings, e.g., surgical needles and catheters, work specific to this domain has been the focus of recent work \cite{ng2025contact,11373832}.Another brand of contact-aided planning is bracing for articulated robots, that is, using contact with the environment to reduce the torque forces on a robot \cite{fang2019exploitation,natarajan2023torque}.


%\cite{yu2024whisker} A survey of whisker sensors fo environment perception.


\subsection{DLO state estimation without vision}
\begin{outline}
  \item Vision is the dominant modality for DLO perception and
    manipulation~\cite{sanchez2018survey}, but it fails under the occlusions we
    target. Force/torque (F/T) and tactile sensing are therefore common in DLO
    insertion and
    routing~\cite{kr-aecsagphi-22,nkot-sditfw-97,slk-ftmdloeffi-22,wjly-cratmp-22,zdpp-dhattfmn-19,lvw-crcrisb-24}.
  \item For an elastic rod in static equilibrium, the shape is determined by
    the wrench at one end~\cite{bretl2014quasistatic}; this enabled real-time,
    vision-free shape estimation of a wire held by two arms from F/T
    sensing~\cite{mishani2022realtime}.
  \item \citet{slk-ftmdloeffi-22} manipulate a DLO by keeping it under tension,
    assuming it is straight when tensioned and that gravity makes free segments
    vertical; F/T readings detect the tension and feed an analytic model for
    estimating the DLO state.
  \item Gap: these works assume a contact-free rod held at both ends, or a DLO
    kept taut. An environment contact that bends the rod adds an unknown
    reaction force that breaks the wrench-to-shape mapping; we add the contact
    location to recover it.
\end{outline}

\subsection{DLO insertion and threading}
\begin{outline}
  \item \citet{nkot-sditfw-97} insert a flexible wire using a stereo camera and
    an F/T sensor, first straightening its free end; the insertion strategy is
    standard, and deformation caused by contact with the environment is not
    discussed.
  \item Tight-tolerance threading of string and rope through
    openings~\cite{wang2015online}; learned DLO insertion with flexibility
    estimated from visual cues~\cite{lc-ldloilfevc-24}.
  \item Tactile insertion: \citet{zdpp-dhattfmn-19} learn to guide a deformed
    DLO through a hole using a 16-taxel tactile sensor, without assuming
    compliance or deforming the DLO. \citet{wjly-cratmp-22} build a library of
    tactile-driven primitives (cable following, pivoting, weaving, insertion)
    using a GelSight sensor, designed for the NIST Assembly Task Board; each
    primitive is a state machine whose transitions are triggered by tactile
    input, and insertion is standard: push until the force drops and a rise in
    torque shows that rotation is impossible. Tactile feedback also enables
    following a cable in the grip~\cite{she2021cable}.
  \item \citet{lvw-crcrisb-24} use reinforcement learning (RL) with F/T sensing
    to push compliant objects, including DLOs, into tunnels of unknown shape; a
    Dynamic Safety Lock computes the forces acting on the end effector to allow
    some pushing but stops when the resistance is too high. This is the only
    work we know of with non-straight passages, but its bends do not exceed
    roughly $\pi/6$, and it addresses the uncertainty of the passage shape
    rather than motion planning or control.
  \item Gap: insertion either relies on vision or pushes along a straight or
    mildly bent passage until the force changes; contact with the passage is
    treated as resistance or uncertainty, not used to redirect the tip, so
    angular threading behind an occluding edge is not addressed.
\end{outline}

\subsection{Exploiting environment contacts}
\begin{outline}
  \item Extrinsic dexterity uses external contacts to extend what a simple
    gripper can do~\cite{dafle2014extrinsic}. For peg-in-hole insertion,
    \citet{kr-aecsagphi-22} estimate the contact line between the peg and the
    hole with a GelSight sensor and feed it to an RL model that selects the
    next action.
  \item For DLOs, environmental contacts and fixtures have been used to shape
    cables~\cite{zhu2020environmental} and to manipulate them blindly with
    force feedback~\cite{slk-ftmdloeffi-22}.
  \item Contacts along slender elastic bodies can be detected and localized
    from kinematics, as done for continuum robots~\cite{bajo2012kinematics}.
  \item Gap: contacts are used to shape or hold the DLO body, or are sensed on a
    rigid peg, not used as a joint for steering a \emph{free tip} to a target
    pose with haptic feedback.
    \todo{Confirm this novelty claim against recent work (2024--2026).}
\end{outline}
